% c05 原书 p.45–58（PDF p.052–065）

\begin{proposition}\label{it:1.3.18}
设 $(X,\cO_X)$ 是概形，$Y$ 是 $X$ 的闭子集。则 $Y$ 上存在唯一的既约概形结构 $(Y,\cO_Y)$，使 $Y$ 成为 $X$ 的闭子概形。若 $(Z,\cO_Z)\to (X,\cO_X)$ 是闭浸入，且它在 $X$ 中的像包含 $Y$，则存在唯一态射 $(Y,\cO_Y)\to (Z,\cO_Z)$，使下图交换：
\begin{center}
\begin{tikzcd}
Y \arrow[r] \arrow[d] & X \\
Z \arrow[ur] &
\end{tikzcd}
\end{center}
我们称此子概形结构为 $Y$ 上的\emph{诱导既约闭子概形结构}\idx{youdaojiyue}{诱导既约概形}。
\end{proposition}

\begin{tnote}
原书 \ref{it:1.3.18} 中唯一态射的源印作 $(Y,\cO_X)$；按上下文应为 $(Y,\cO_Y)$，译文按后者排印。[依据：原书 p.45]
\end{tnote}

\begin{proof}
先考虑 $X=\Spec A$ 是仿射的情形。则 $Y=V(\fa)$，其中 $\fa$ 是 $A$ 的满足 $\fa=\sqrt{\fa}$ 的唯一理想。注意 $\Spec A/\fa$ 是 $\Spec A$ 的既约闭子概形，其底拓扑空间为 $Y$。这证明了诱导既约闭子概形结构的存在性。由 \ref{it:1.3.16}(ii)，底拓扑空间为 $Y$ 的 $X$ 的任意闭子概形都具有形式 $\Spec A/\fa$，其中 $\fa$ 是满足 $Y=V(\fa)$ 的某个理想。若该闭子概形既约，则必有 $\fa=\sqrt{\fa}$。这样的 $\fa$ 是唯一的。这证明了诱导既约闭子概形结构的唯一性。

现考虑一般情形。用仿射开子概形 $U_i=\Spec A_i$（$i\in I$）覆盖 $X$。我们已知，对每个 $i$，仿射概形 $U_i$ 的闭子集 $Y\cap U_i$ 有诱导既约闭子概形结构 $(Y\cap U_i,\cO_{Y\cap U_i})$。我们断言：对任意 $i,j$，闭浸入
\[
(Y\cap U_i\cap U_j,\cO_{Y\cap U_i}|_{Y\cap U_i\cap U_j})\to (U_i\cap U_j,\cO_{U_i\cap U_j})
\]
与闭浸入
\[
(Y\cap U_i\cap U_j,\cO_{Y\cap U_j}|_{Y\cap U_i\cap U_j})\to (U_i\cap U_j,\cO_{U_i\cap U_j})
\]
之间存在唯一同构。于是可以把闭子概形 $(Y\cap U_i,\cO_{Y\cap U_i})$（$i\in I$）粘起来，得到 $Y$ 上的既约闭子概形结构。（参看下面的 \ref{it:1.3.19}。）为证该断言，用仿射开子概形 $U_{ijk}$（$k\in K_{ij}$）覆盖 $U_i\cap U_j$。对任意 $i,j,k$，$(Y\cap U_{ijk},\cO_{Y\cap U_i}|_{Y\cap U_{ijk}})$ 与 $(Y\cap U_{ijk},\cO_{Y\cap U_j}|_{Y\cap U_{ijk}})$ 都是既约的，且它们的底拓扑空间都是仿射概形 $U_{ijk}$ 的闭子集 $Y\cap U_{ijk}$。由仿射情形下诱导既约闭子概形结构的唯一性，闭浸入
\[
(Y\cap U_{ijk},\cO_{Y\cap U_i}|_{Y\cap U_{ijk}})\to (U_{ijk},\cO_{U_{ijk}})
\]
与闭浸入
\[
(Y\cap U_{ijk},\cO_{Y\cap U_j}|_{Y\cap U_{ijk}})\to (U_{ijk},\cO_{U_{ijk}})
\]
之间存在唯一同构。因 $U_{ijk}$（$k\in K_{ij}$）覆盖 $U_i\cap U_j$，我们的断言成立。

诱导既约闭子概形结构的唯一性也可归结为仿射情形来证明。细节留给读者。
\end{proof}

在上述证明中，我们通过粘合子概形结构得到了诱导既约闭子概形结构。下面把粘合过程概括为一个命题，其证明留给读者：

\begin{proposition}\label{it:1.3.19}
(i) 设 $U_i$（$i\in I$）是一族拓扑空间。假设对任意一对 $i\ne j$，给定 $U_i$ 的开子集 $U_{ij}$ 与同胚 $\varphi_{ij}\colon U_{ij}\to U_{ji}$，满足以下性质：
\begin{enumerate}[label=(\alph*),leftmargin=2.8em]
  \item 对任意 $i\ne j$，有 $\varphi_{ij}=\varphi_{ji}^{-1}$。
  \item 对任意互不相同的 $i,j,k$，有 $\varphi_{ij}(U_{ij}\cap U_{ik})=U_{ji}\cap U_{jk}$，且在 $U_{ij}\cap U_{ik}$ 上 $\varphi_{jk}\varphi_{ij}=\varphi_{ik}$。
\end{enumerate}
则存在拓扑空间 $X$ 与嵌入 $\varphi_i\colon U_i\to X$，使 $\varphi_i(U_i)$ 是开集并构成 $X$ 的覆盖，$\varphi_i(U_{ij})=\varphi_i(U_i)\cap\varphi_j(U_j)$，且在 $U_{ij}$ 上 $\varphi_j\varphi_{ij}=\varphi_i$。$X$ 在同胚意义下唯一。我们称 $\{\varphi_{ij}\}$ 为粘合数据，称 $X$ 由沿 $\{U_{ij}\}$ 通过 $\{\varphi_{ij}\}$ 粘合 $\{U_i\}$ 得到。

(ii) 设 $X$ 是拓扑空间，$\{U_i\}_{i\in I}$ 是 $X$ 的开覆盖。假设对每个 $i$，给定 $U_i$ 上的层 $\cF_i$；对每对 $i,j$，给定层同构 $\phi_{ij}\colon \cF_i|_{U_i\cap U_j}\to \cF_j|_{U_i\cap U_j}$，满足以下性质：
\begin{enumerate}[label=(\alph*),leftmargin=2.8em]
  \item 对每个 $i$，有 $\phi_{ii}=\id_{\cF_i}$。
  \item 对任意三元组 $i,j,k$，在 $\cF_i|_{U_i\cap U_j\cap U_k}$ 上有 $\phi_{ik}=\phi_{jk}\phi_{ij}$。
\end{enumerate}
则存在 $X$ 上的层 $\cF$ 与同构 $\phi_i\colon \cF|_{U_i}\to \cF_i$，使在 $\cF|_{U_i\cap U_j}$ 上有 $\phi_j=\phi_{ij}\phi_i$。$\cF$ 在同构意义下唯一。我们称 $\{\phi_{ij}\}$ 为粘合数据，称 $\cF$ 由通过 $\{\phi_{ij}\}$ 粘合 $\{\cF_i\}$ 得到。

(iii) 设 $U_i$（$i\in I$）是一族概形。假设对任意一对 $i\ne j$，给定 $U_i$ 的开子概形 $U_{ij}$ 与概形同构 $\varphi_{ij}\colon U_{ij}\to U_{ji}$，满足以下性质：
\begin{enumerate}[label=(\alph*),leftmargin=2.8em]
  \item 对任意 $i\ne j$，有 $\varphi_{ij}=\varphi_{ji}^{-1}$。
  \item 对任意互不相同的 $i,j,k$，有 $\varphi_{ij}(U_{ij}\cap U_{ik})=U_{ji}\cap U_{jk}$，且在 $U_{ij}\cap U_{ik}$ 上 $\varphi_{jk}\varphi_{ij}=\varphi_{ik}$。
\end{enumerate}
则存在概形 $X$ 与开浸入 $\varphi_i\colon U_i\to X$，使 $\varphi_i(U_i)$ 构成 $X$ 的覆盖，$\varphi_i(U_{ij})=\varphi_i(U_i)\cap\varphi_j(U_j)$，且在 $U_{ij}$ 上 $\varphi_j\varphi_{ij}=\varphi_i$。$X$ 在同构意义下唯一。我们称 $\{\varphi_{ij}\}$ 为粘合数据，称 $X$ 由沿 $\{U_{ij}\}$ 通过 $\{\varphi_{ij}\}$ 粘合 $\{U_i\}$ 得到。
\end{proposition}

设 $S$ 是概形。\emph{$S$-概形}或\emph{$S$ 上的概形}\idx{sgaixing}{S-概形}是指概形 $X$ 连同态射 $X\to S$。从 $S$-概形 $X$ 到 $S$-概形 $Y$ 的\emph{$S$-态射}\idx{staishe}{S-态射}是指使图表
\begin{center}
\begin{tikzcd}
X \arrow[r] \arrow[dr] & Y \arrow[d] \\
& S
\end{tikzcd}
\end{center}
交换的态射 $X\to Y$。用显然的方式定义 $S$-态射的复合。于是得到 $S$-概形范畴。设 $X$ 与 $Y$ 是 $S$-概形。它们在 $S$-概形范畴中的积称为 $X$ 与 $Y$ 在 $S$ 上的\emph{纤维积}\idx{xianweiji}{纤维积}。它是 $S$-概形 $X\times_S Y$，连同两个称为\emph{投影}\idx{touying}{投影}的 $S$-态射 $p\colon X\times_S Y\to X$ 与 $q\colon X\times_S Y\to Y$，具有以下泛性质：对任意 $S$-概形 $Z$ 与任意 $S$-态射 $f\colon Z\to X$ 及 $g\colon Z\to Y$，存在唯一 $S$-态射 $h\colon Z\to X\times_S Y$，使 $ph=f$ 且 $qh=g$。有交换图表
\begin{center}
\begin{tikzcd}
& X\times_S Y \arrow[dl, "p"'] \arrow[dr, "q"] & \\
X \arrow[dr] & & Y \arrow[dl] \\
& S &
\end{tikzcd}
\end{center}
若一个图表具有上述形式，则称它为\emph{笛卡尔图}\idx{dikaertu}{笛卡尔图}。

注意由 \ref{it:1.2.8}，对任意概形 $X$，有唯一态射 $X\to \Spec\bZ$。故概形范畴与 $\Spec\bZ$-概形范畴重合。

两个概形 $X$ 与 $Y$ 在概形范畴中的积记为 $X\times Y$。它与 $X\times_{\Spec\bZ}Y$ 相同。若 $S=\Spec R$ 是仿射概形，我们也把 $S$-概形（resp.\ $S$-态射）称为 $R$-概形（resp.\ $R$-态射），并常把 $X\times_S Y$ 写成 $X\times_R Y$。

\begin{proposition}\label{it:1.3.20}
对任意 $S$-概形 $X$ 与 $Y$，它们在 $S$ 上的纤维积存在，且在同构意义下唯一。
\end{proposition}

\begin{proof}
唯一性由纤维积的泛性质得到。

\emph{第 1 步}。设 $X=\Spec A$，$Y=\Spec B$，$S=\Spec C$ 都是仿射的。则由 \ref{it:1.2.8} 可验证 $X\times_S Y=\Spec(A\tensor_C B)$，且 $p$ 与 $q$ 分别由典范 $C$-代数同态 $A\to A\tensor_C B$ 与 $B\to A\tensor_C B$ 诱导。

\emph{第 2 步}。设 $X$ 与 $Y$ 是使 $X\times_S Y$ 存在的 $S$-概形。则对 $X$ 的任意开子概形 $U$，$p^{-1}(U)$ 是 $U$ 与 $Y$ 在 $S$ 上的纤维积。

\emph{第 3 步}。设 $X$ 可被开子概形 $U_i$（$i\in I$）覆盖，且每个 $U_i\times_S Y$ 存在。则 $X\times_S Y$ 存在。事实上，由纤维积的唯一性与第 2 步，对任意 $i,j\in I$，可把 $(U_i\cap U_j)\times_S Y$ 视为 $U_i\times_S Y$ 与 $U_j\times_S Y$ 的开子概形。于是可沿 $\{U_{ij}\times_S Y\}$ 粘合 $\{U_i\times_S Y\}$。然后可以验证所得概形就是 $X\times_S Y$。

\emph{第 4 步}。用仿射开子概形覆盖 $X$，并应用第 1、3 步，可知若 $S$ 与 $Y$ 是仿射的，则 $X\times_S Y$ 存在。

\emph{第 5 步}。用仿射开子概形覆盖 $Y$，并对 $X$ 与 $Y$ 互换后应用第 4 步与第 3 步，可知若 $S$ 是仿射的，则 $X\times_S Y$ 存在。

\emph{第 6 步}。一般地，用仿射开子概形 $S_i$（$i\in I$）覆盖 $S$，并设 $X_i$ 与 $Y_i$ 分别是 $S_i$ 在 $X$ 与 $Y$ 中的原像。则由第 5 步，每个 $X_i\times_{S_i}Y_i$ 存在。可以证明 $X_i\times_S Y\cong X_i\times_{S_i}Y_i$。应用第 3 步，即知 $X\times_S Y$ 存在。
\end{proof}

设给定概形态射 $S'\to S$。对任意 $S$-概形 $X$，投影 $X\times_S S'\to S'$ 在 $X\times_S S'$ 上定义了 $S'$-概形结构。对 $S$-概形之间的任意 $S$-态射 $f\colon X\to Y$，设 $f\times\id_{S'}\colon X\times_S S'\to Y\times_S S'$ 为使下图交换的唯一 $S'$-态射：
\begin{center}
\begin{tikzcd}
X\times_S S' \arrow[r] \arrow[d] & Y\times_S S' \arrow[d] \\
X \arrow[r, "f"] & Y
\end{tikzcd}
\end{center}
于是得到从 $S$-概形范畴到 $S'$-概形范畴的函子。此函子称为从 $S$ 到 $S'$ 的\emph{基变换}\idx{jibianhuan}{基变换}，或\emph{基扩张}。

\begin{proposition}\label{it:1.3.21}
设 $f\colon X\to Y$ 是态射，$y$ 是 $Y$ 中的点，$k(y)$ 是局部环 $\cO_{Y,y}$ 的剩余域。则投影
\[
X\times_Y \Spec k(y)\to X
\]
在底拓扑空间上诱导 $X\times_Y \Spec k(y)$ 与 $f^{-1}(y)$ 的同胚。
\end{proposition}

我们称 $k(y)$-概形 $X\times_Y \Spec k(y)$ 为 $f$ 在 $y$ 上的\emph{纤维}\idx{xianwei}{纤维}。

\begin{proof}
问题具有局部性。可设 $X=\Spec B$ 与 $Y=\Spec A$ 是仿射的，且 $f\colon X\to Y$ 由同态 $\phi\colon A\to B$ 诱导。设 $\fp$ 是 $A$ 中对应于 $y$ 的素理想。则有
\[
X\times_Y \Spec k(y)\cong \Spec(B\tensor_A A_\fp/\fp A_\fp)\cong \Spec B_\fp/\fp B_\fp.
\]
有以下一一对应：
\[
\begin{aligned}
&\{\fp \text{ 中的 } B_\fp/\fp B_\fp \text{ 的素理想}\}\\
&\cong \{B_\fp \text{ 的包含 } \fp B_\fp \text{ 的素理想}\}\\
&\cong \{\fq \text{ 是 } B \text{ 的素理想} \mid \fq\cap\phi(A-\fp)=\emptyset \text{ 且 } \fq\supset\phi(\fp)\}\\
&\cong \{\fq \text{ 是 } B \text{ 的素理想} \mid \phi^{-1}(\fq)=\fp\}.
\end{aligned}
\]
这说明
\[
X\times_Y \Spec k(y)=\Spec B_\fp/\fp B_\fp\to X=\Spec B
\]
在 $X\times_Y \Spec k(y)$ 与 $f^{-1}(y)$ 之间诱导一一对应。为说明它是嵌入，把它分解为复合
\[
\Spec B_\fp/\fp B_\fp\to \Spec B_\fp\to \Spec B.
\]
注意第一个箭头是闭浸入，第二个箭头由 \ref{it:1.2.2} 后 1.2 节习题 2 是底拓扑空间上的嵌入。
\end{proof}

设 $X$ 是拓扑空间。则 $X$ 是 Hausdorff 空间，当且仅当对角集 $\{(x,x)\mid x\in X\}$ 在 $X\times X$ 中闭。分离态射的概念对应于 Hausdorff 拓扑空间的概念。

设 $f\colon X\to Y$ 是概形的态射。定义\emph{对角态射}\idx{duijiao}{对角态射}
\[
\Delta_{X/Y}\colon X\to X\times_Y X
\]
为使
\[
p\Delta_{X/Y}=q\Delta_{X/Y}=\id_X
\]
的唯一态射。常把 $\Delta_{X/Y}$ 记为 $\Delta$。若 $\Delta_{X/Y}$ 是闭浸入，则称 $f\colon X\to Y$ 为\emph{分离态射}\idx{fenli}{分离态射}，或说 $X$ 在 $Y$ 上\emph{分离}。若典范态射 $X\to \Spec\bZ$ 分离，则称概形 $X$ 为\emph{分离概形}。

\begin{proposition}\label{it:1.3.22}
设 $f\colon \Spec B\to \Spec A$ 是仿射概形之间的态射。则 $f$ 分离。
\end{proposition}

\begin{proof}
不难看出，$\Delta\colon \Spec B\to \Spec(B\tensor_A B)$ 对应于同态
\[
B\tensor_A B\to B,\quad b_1\tensor b_2\mapsto b_1b_2,
\]
它是满射。故由 \ref{it:1.3.16}(i)，$\Delta$ 是闭浸入。
\end{proof}

\begin{proposition}\label{it:1.3.23}
设 $f\colon X\to Y$ 是概形的态射。
\begin{enumerate}[label=(\roman*),leftmargin=2.8em]
  \item 对角态射 $\Delta\colon X\to X\times_Y X$ 是浸入。
  \item $f\colon X\to Y$ 分离，当且仅当 $\Delta(X)$ 是 $X\times_Y X$ 的闭子集。
\end{enumerate}
\end{proposition}

\begin{proof}
(i) 用仿射开子集 $V_i$（$i\in I$）覆盖 $Y$，并用仿射开子集 $U_{ij}$（$j\in J_i$）覆盖每个 $f^{-1}(V_i)$。设 $U$ 是 $p^{-1}(U_{ij})\cap q^{-1}(U_{ij})$（$i\in I$，$j\in J_i$）的并。则 $\Delta(X)$ 包含在 $U$ 中。故可把 $\Delta$ 分解为复合
\[
X\to U\to X\times_Y X.
\]
第二个箭头是开浸入。为证 $\Delta$ 是浸入，只需证明第一个箭头是闭浸入。只需证明每个
\[
\Delta^{-1}(p^{-1}(U_{ij})\cap q^{-1}(U_{ij}))\xrightarrow{\Delta}p^{-1}(U_{ij})\cap q^{-1}(U_{ij})
\]
是闭浸入。我们有交换图表
\begin{center}
\begin{tikzcd}
\Delta^{-1}(p^{-1}(U_{ij})\cap q^{-1}(U_{ij})) \arrow[r, "\Delta"] \arrow[d, equal] & p^{-1}(U_{ij})\cap q^{-1}(U_{ij}) \arrow[d, "\cong"] \\
U_{ij} \arrow[r, "\Delta_{U_{ij}/V_i}"] & U_{ij}\times_{V_i} U_{ij}.
\end{tikzcd}
\end{center}
由 \ref{it:1.3.22}，下水平箭头是闭浸入。于是我们的断言成立。

(ii) 由 (i)，只需证明浸入 $f\colon X\to Y$ 是闭浸入，当且仅当 $f(X)$ 是 $Y$ 的闭子集。「仅当」部分显然。为证「当」部分，把 $f\colon X\to Y$ 分解为复合 $X\to V\to Y$，使第一个箭头是闭浸入，第二个箭头是开浸入。因 $f(X)$ 闭，$Y$ 有开覆盖 $\{V,Y-f(X)\}$。注意 $f^{-1}(V)\xrightarrow{f}V$ 与 $f^{-1}(Y-f(X))\xrightarrow{f}Y-f(X)$ 是闭浸入，因为第一个态射与 $X\to V$ 重合，且 $f^{-1}(Y-f(X))=\emptyset$。故 $f$ 是闭浸入。
\end{proof}

下面命题的证明留给读者：

\begin{proposition}\label{it:1.3.24}
设
\begin{center}
\begin{tikzcd}
X \arrow[dr] & & Y \arrow[dl] \\
& T \arrow[d] & \\
& S &
\end{tikzcd}
\end{center}
是态射。则图表
\begin{center}
\begin{tikzcd}
X\times_T Y \arrow[r] \arrow[d] & X\times_S Y \arrow[d] \\
T \arrow[r, "\Delta"] & T\times_S T
\end{tikzcd}
\end{center}
是笛卡尔图，其中上水平箭头是使它与 $X\times_S Y$ 到两个因子的两个投影复合后，分别等于 $X\times_T Y$ 到两个因子的两个投影的唯一态射；右竖箭头是使它与 $T\times_S T$ 到两个因子的两个投影复合后，分别等于态射 $X\to T$ 与 $Y\to T$ 的唯一态射。
\end{proposition}
设 $X$ 与 $Y$ 是 $S$-概形，$f\colon X\to Y$ 是 $S$-态射。把上述命题应用于
\begin{center}
\begin{tikzcd}
X \arrow[dr, "f"'] & & Y \arrow[dl, "\id"] \\
& Y \arrow[d] & \\
& S &
\end{tikzcd}
\end{center}
并注意 $X\times_Y Y=X$，我们得到笛卡尔图
\begin{center}
\begin{tikzcd}
X \arrow[r] \arrow[d] & X\times_S Y \arrow[d] \\
Y \arrow[r, "\Delta"] & Y\times_S Y.
\end{tikzcd}
\end{center}
上水平箭头称为 $S$-态射 $f\colon X\to Y$ 的\emph{图像}\idx{tuxiang}{图像}，记为 $\Gamma_f\colon X\to X\times_S Y$。它是使
\[
p\Gamma_f=\id_X,\quad q\Gamma_f=f
\]
的唯一态射。它由 $\Delta_{Y/S}$ 经基变换 $X\times_S Y\to Y\times_S Y$ 得到。利用 \ref{it:1.3.16}(ii)，可以证明浸入（resp.\ 开浸入，resp.\ 闭浸入）的基变换仍是浸入（resp.\ 开浸入，resp.\ 闭浸入）。故由 \ref{it:1.3.23}(i)，$\Gamma_f$ 是浸入。当 $Y$ 在 $S$ 上分离时，$\Gamma_f$ 是闭浸入。

若对角态射 $\Delta\colon X\to X\times_Y X$ 拟紧，则称概形的态射 $f\colon X\to Y$ 为\emph{拟分离}\idx{nifenli}{拟分离}的。若典范态射 $X\to \Spec\bZ$ 拟分离，则称概形 $X$ 为\emph{拟分离概形}。

\begin{proposition}\label{it:1.3.25}
设 $f\colon X\to \Spec A$ 是从 $X$ 到仿射概形的态射，$U$ 与 $V$ 是 $X$ 的两个仿射开子概形。
\begin{enumerate}[label=(\roman*),leftmargin=2.8em]
  \item 若 $f$ 分离，则 $U\cap V$ 是仿射的。$A$-代数 $\cO_X(U\cap V)$ 由限制 $\cO_X(U)\to \cO_X(U\cap V)$ 与 $\cO_X(V)\to \cO_X(U\cap V)$ 的像生成。
  \item 若 $f$ 拟分离，则 $U\cap V$ 可被有限多个仿射开子概形覆盖。
\end{enumerate}
\end{proposition}

\begin{proof}
(i) 因对角态射 $\Delta\colon X\to X\times_{\Spec A}X$ 是闭浸入，我们有闭浸入
\[
U\cap V=\Delta^{-1}(p^{-1}(U)\cap q^{-1}(V))\xrightarrow{\Delta}p^{-1}(U)\cap q^{-1}(V)\cong U\times_{\Spec A}V.
\]
但 $U\times_{\Spec A}V$ 是仿射的，故由 \ref{it:1.3.16}(ii)，$U\cap V$ 也是仿射的，且同态 $\cO_{U\times_{\Spec A}V}(U\times_{\Spec A}V)\to \cO_X(U\cap V)$ 是满射。此满射可与由 $A$-代数同态 $\cO_X(U)\to \cO_X(U\cap V)$ 与 $\cO_X(V)\to \cO_X(U\cap V)$ 诱导的同态
\[
\cO_X(U)\tensor_A \cO_X(V)\to \cO_X(U\cap V)
\]
等同。故 $\cO_X(U\cap V)$ 由它们的像生成。

(ii) 注意 $p^{-1}(U)\cap q^{-1}(V)\cong U\times_{\Spec A}V$ 是 $X\times_{\Spec A}X$ 的仿射开子概形。因 $\Delta\colon X\to X\times_{\Spec A}X$ 拟紧，$U\cap V=\Delta^{-1}(p^{-1}(U)\cap q^{-1}(V))$ 拟紧，从而可被有限多个仿射开子概形覆盖。
\end{proof}

若概形的态射 $f\colon X\to Y$ 满足以下性质，则称 $f$ 为\emph{固有时态射}\idx{guyou}{固有态射}，或说 $X$ 在 $Y$ 上\emph{固有}：
\begin{enumerate}[label=(\roman*),leftmargin=2.8em]
  \item $f$ 有限型。
  \item $f$ 分离。
  \item 对任意态射 $Y'\to Y$，$f$ 的基变换 $f'\colon X\times_Y Y'\to Y'$ 是底拓扑空间上的闭映射。若 $f$ 具有此性质，则称 $f$ \emph{泛闭}\idx{fanbi}{泛闭}。
\end{enumerate}

\begin{proposition}\label{it:1.3.26}
\begin{enumerate}[label=(\roman*),leftmargin=2.8em]
  \item 态射 $f\colon X\to Y$ 分离（resp.\ 固有），当且仅当存在 $Y$ 的开覆盖 $\{V_i\}_{i\in I}$，使每个 $f^{-1}(V_i)\to V_i$ 分离（resp.\ 固有）。
  \item 浸入分离。闭浸入固有。
  \item 两个分离态射（resp.\ 固有态射）的复合仍分离（resp.\ 固有）。
  \item 设 $f\colon X\to Y$ 与 $Y'\to Y$ 是态射，$f'\colon X\times_Y Y'\to Y'$ 是 $f$ 的基变换。若 $f$ 分离（resp.\ 固有），则 $f'$ 分离（resp.\ 固有）。
  \begin{center}
  \begin{tikzcd}
  X\times_Y Y' \arrow[r, "f'"] \arrow[d] & Y' \arrow[d] \\
  X \arrow[r, "f"] & Y.
  \end{tikzcd}
  \end{center}
  \item 设 $f\colon X\to Y$ 与 $f'\colon X'\to Y'$ 是 $S$-概形之间的两个分离（resp.\ 固有）$S$-态射。则 $f\times f'\colon X\times_S X'\to Y\times_S Y'$ 分离（resp.\ 固有），其中 $f\times f'$ 是使以下图表交换的唯一态射：
  \begin{center}
  \begin{tikzcd}
  X\times_S Y \arrow[r, "f\times f'"] \arrow[d] & X'\times_S Y' \arrow[d] \\
  X \arrow[r, "f"] & X',
  \end{tikzcd}
  \qquad
  \begin{tikzcd}
  X\times_S Y \arrow[r, "f\times f'"] \arrow[d] & X'\times_S Y' \arrow[d] \\
  Y \arrow[r, "f'"] & Y'.
  \end{tikzcd}
  \end{center}
  \item 设 $f\colon X\to Y$ 与 $g\colon Y\to Z$ 是两个态射。若 $gf$ 分离，则 $f$ 分离。若 $gf$ 固有且 $g$ 分离，则 $f$ 固有。
\end{enumerate}
\end{proposition}

\begin{tnote}
原书 \ref{it:1.3.26}(v) 的两个图表左端均印作 $X\times_S Y$；按 $f\times f'$ 的靶为 $Y\times_S Y'$ 以及下方两行动机判断，左端应是 $X\times_S X'$。此处照原书排印，特此说明。[依据：原书 p.52]
\end{tnote}

\begin{proof}
(i) 以及关于固有态射的断言，其证明留给读者。

(ii) 设 $X\to Y$ 是浸入。利用它是概形范畴中的单态射这一事实，可以证明 $X\times_Y X\cong X$，且 $\Delta\colon X\to X\times_Y X$ 可等同于 $\id_X$。于是 $\Delta$ 是闭浸入，$X$ 在 $Y$ 上分离。

(iii) 设 $X\to Y$ 与 $Y\to Z$ 是两个态射。由 \ref{it:1.3.24}，有笛卡尔图
\begin{center}
\begin{tikzcd}
X\times_Y X \arrow[r] \arrow[d] & X\times_Z X \arrow[d] \\
Y \arrow[r, "\Delta_{Y/Z}"] & Y\times_Z Y.
\end{tikzcd}
\end{center}
若 $Y\to Z$ 分离，则 $\Delta_{Y/Z}$ 是闭浸入，从而其基变换 $X\times_Y X\to X\times_Z X$ 是闭浸入。有交换图表
\begin{center}
\begin{tikzcd}
X \arrow[r, "\Delta_{X/Y}"] \arrow[dr, "\Delta_{X/Z}"'] & X\times_Y X \arrow[d] \\
& X\times_Z X.
\end{tikzcd}
\end{center}
若 $X\to Y$ 也分离，则 $\Delta_{X/Z}$ 是闭浸入，因为它是闭浸入 $\Delta_{X/Y}$ 与 $X\times_Y X\to X\times_Z X$ 的复合。于是 $X\to Z$ 分离。

(iv) 设 $X'=X\times_Y Y'$。有笛卡尔图
\begin{center}
\begin{tikzcd}
X' \arrow[r, "\Delta_{X'/Y'}"] \arrow[d] & X'\times_{Y'}X' \arrow[d] \\
X \arrow[r, "\Delta_{X/Y}"] & X\times_Y X.
\end{tikzcd}
\end{center}
故若 $\Delta_{X/Y}$ 是闭浸入，则 $\Delta_{X'/Y'}$ 是闭浸入。

(v) 设 $f$ 与 $f'$ 分离。有
\[
f\times f'=(f\times\id_{Y'})(\id_X\times f').
\]
为证 $f\times f'$ 分离，由 (iii) 只需证明 $f\times\id_{Y'}$ 与 $\id_X\times f'$ 分离。不难验证
\begin{center}
\begin{tikzcd}
X\times_S Y' \arrow[r] \arrow[d] & X \arrow[d, "f"] \\
Y\times_S Y' \arrow[r] & Y
\end{tikzcd}
\end{center}
是笛卡尔图（即 $X\times_Y(Y\times_S Y')=X\times_S Y'$）。故 $f\times\id_{Y'}$ 是 $f$ 的基变换，从而分离。类似地，$\id_X\times f'$ 是 $f'$ 的基变换，从而分离。

(vi) 有 $f=q\Gamma_f$，其中 $\Gamma_f\colon X\to X\times_Z Y$ 是 $Z$-态射 $f$ 的图像，$q\colon X\times_Z Y\to Y$ 是投影。注意 $q$ 是 $gf$ 的基变换，从而分离。$\Gamma_f$ 是浸入，从而分离。故它们的复合 $f$ 分离。
\end{proof}

考虑环 $\bZ[x_0,\dots,x_n]$，其分次由多项式的次数定义。把 $\Proj \bZ[x_0,\dots,x_n]$ 记为 $\bP^n_{\bZ}$。对任意概形 $Y$，定义 $Y$ 上的\emph{射影空间}\idx{sheyingkongjian}{射影空间}为 $Y$-概形 $\bP^n_Y=\bP^n_{\bZ}\times Y$。若态射 $f\colon X\to Y$ 可分解为复合
\[
X\to \bP^n_Y\to Y,
\]
其中 $X\to \bP^n_Y$ 是闭浸入，$\bP^n_Y\to Y$ 是投影，则称 $f$ \emph{射影}\idx{sheying}{射影态射}，或说 $X$ 在 $Y$ 上射影。若可作上述分解且 $X\to \bP^n_Y$ 是浸入，则称 $f$ \emph{拟射影}\idx{nisheying}{拟射影}。

\begin{proposition}\label{it:1.3.27}
射影态射固有。
\end{proposition}

\begin{proof}
考虑到 \ref{it:1.3.26}(ii),(iii),(iv)，只需证明 $\bP^n_{\bZ}\to \Spec\bZ$ 有限型、分离且泛闭。

用 $D_+(x_i)\cong \Spec \bZ[x_0,\dots,x_n]_{(x_i)}$（$i=0,\dots,n$）覆盖 $\bP^n_{\bZ}$。有
\[
\bZ[x_0,\dots,x_n]_{(x_i)}\cong \bZ\Bigl[\frac{x_0}{x_i},\dots,\frac{x_n}{x_i}\Bigr].
\]
故 $\bZ[x_0,\dots,x_n]_{(x_i)}$ 是有限生成的 $\bZ$-代数。于是 $\bP^n_{\bZ}\to \Spec\bZ$ 有限型。

用 $p^{-1}(D_+(x_i))\cap q^{-1}(D_+(x_j))$（$i,j=0,\dots,n$）覆盖 $\bP^n_{\bZ}\times\bP^n_{\bZ}$。有
\[
\begin{aligned}
\Delta^{-1}(p^{-1}(D_+(x_i))\cap q^{-1}(D_+(x_j)))
&=D_+(x_i)\cap D_+(x_j)=D_+(x_ix_j)\\
&\cong \Spec \bZ[x_0,\dots,x_n]_{(x_ix_j)},
\end{aligned}
\]
以及
\[
\begin{aligned}
p^{-1}(D_+(x_i))\cap q^{-1}(D_+(x_j))
&\cong D_+(x_i)\times D_+(x_j)\\
&\cong \Spec(\bZ[x_0,\dots,x_n]_{(x_i)}\tensor_{\bZ}\bZ[x_0,\dots,x_n]_{(x_j)}).
\end{aligned}
\]
由于典范同态
\[
\bZ[x_0,\dots,x_n]_{(x_i)}\tensor_{\bZ}\bZ[x_0,\dots,x_n]_{(x_j)}\to \bZ[x_0,\dots,x_n]_{(x_ix_j)},\quad
\frac{x_k}{x_i}\tensor\frac{x_l}{x_j}\mapsto \frac{x_kx_l}{x_ix_j}
\]
是满射，态射
\[
\Delta\colon \Delta^{-1}(p^{-1}(D_+(x_i))\cap q^{-1}(D_+(x_j)))\to p^{-1}(D_+(x_i))\cap q^{-1}(D_+(x_j))
\]
是闭浸入。故 $\Delta\colon \bP^n_{\bZ}\to \bP^n_{\bZ}\times\bP^n_{\bZ}$ 是闭浸入。于是 $\bP^n_{\bZ}\to \Spec\bZ$ 分离。

为证 $\bP^n_{\bZ}\to \Spec\bZ$ 泛闭，只需证明对任意环 $A$，典范态射 $\Proj A[x_0,\dots,x_n]\to \Spec A$ 是底拓扑空间上的闭映射。对 $A[x_0,\dots,x_n]$ 的任意齐次理想 $\fa$，为证 $V_+(\fa)$ 的像在 $\Spec A$ 中闭，只需证明
\[
\Proj A[x_0,\dots,x_n]/\fa\to \Spec A
\]
的像闭。设 $S=A[x_0,\dots,x_n]/\fa$，并设 $\bar x_i$ 是 $x_i$ 在 $S$ 中的像。把典范态射 $\Proj S\to \Spec A$ 记为 $f$。设 $\fp\in \Spec A$ 不在 $f$ 的像中。则 $f^{-1}(\fp)\cap D_+(\bar x_i)$ 为空。有
\[
f^{-1}(\fp)\cap D_+(\bar x_i)\cong \Spec(S_{(\bar x_i)}\tensor_A A_\fp/\fp A_\fp).
\]
由 \ref{it:1.2.2}(i)，在 $S_{(\bar x_i)}\tensor_A A_\fp/\fp A_\fp$ 中有 $0=1$，从而在
\[
S_{\bar x_i}\tensor_A A_\fp/\fp A_\fp\cong (S\tensor_A A_\fp/\fp A_\fp)_{\bar x_i}
\]
中也有 $0=1$。故对每个 $i$，$\bar x_i\tensor 1$ 在 $S\tensor_A A_\fp/\fp A_\fp$ 中幂零。因此可找到自然数 $N$，使对每个 $k\ge N$ 有 $S_k\tensor_A A_\fp/\fp A_\fp=0$。由中山引理，对每个 $k\ge N$ 有 $S_k\tensor_A A_\fp=0$，从而由下面的 \ref{it:1.3.28} 有 $\mathrm{Ann}_A(S_k)\not\subset \fp$。反之，若存在自然数 $N$，使对每个 $k\ge N$ 有 $\mathrm{Ann}_A(S_k)\not\subset \fp$，则把上述论证倒过来，可证 $\fp$ 不在 $f$ 的像中。故 $\fp\notin \operatorname{im}f$，当且仅当存在自然数 $N$，使对每个 $k\ge N$ 有 $\mathrm{Ann}_A(S_k)\not\subset \fp$。显然 $\mathrm{Ann}_A(S_k)\subset \mathrm{Ann}_A(S_{k+1})$。于是有以下等价：
\[
\begin{aligned}
&\fp\in \operatorname{im}f\\
\Leftrightarrow{}&\text{对每个自然数 } N \text{，存在 } k\ge N \text{ 使 } \mathrm{Ann}_A(S_k)\subset \fp\\
\Leftrightarrow{}&\text{对每个自然数 } N \text{，有 } \mathrm{Ann}_A(S_N)\subset \fp\\
\Leftrightarrow{}&\bigcup_N \mathrm{Ann}_A(S_N)\subset \fp.
\end{aligned}
\]
故 $\operatorname{im}f=V(\bigcup_N \mathrm{Ann}_A(S_N))$，它是闭集。
\end{proof}

\begin{lemma}\label{it:1.3.28}
对任意有限生成的 $A$-模 $M$，有
\[
\{\fp\in \Spec A\mid M_\fp=0\}=\Spec A-V(\mathrm{Ann}_A(M)),
\]
其中 $\mathrm{Ann}_A(M)=\{r\in A\mid rx=0,\ \forall x\in M\}$ 是 $M$ 的零化子。
\end{lemma}

\begin{proof}
设 $x_1,\dots,x_n$ 是 $M$ 的有限生成元族。若 $M_\fp=0$，则对每个 $i$ 有 $\frac{x_i}{1}=0$ 在 $M_\fp$ 中成立，从而存在 $s_i\notin \fp$ 使 $s_ix_i=0$。设 $s=s_1\cdots s_n$。则 $s\notin \fp$，且对每个 $i$ 有 $sx_i=0$，从而 $s\in \mathrm{Ann}_A(M)$。故 $\mathrm{Ann}_A(M)\not\subset \fp$。

反之，若 $\mathrm{Ann}_A(M)\not\subset \fp$，则可选取 $s\in \mathrm{Ann}_A(M)$ 但 $s\notin \fp$。对任意 $\frac{x}{t}\in M_\fp$，有 $\frac{x}{t}=\frac{sx}{st}=0$。故 $M_\fp=0$。
\end{proof}
\begin{proposition}\label{it:1.3.29}
\begin{enumerate}[label=(\roman*),leftmargin=2.8em]
  \item 闭浸入射影。
  \item 射影态射的复合射影。
  \item 设 $f\colon X\to Y$ 与 $Y'\to Y$ 是态射，$f'\colon X\times_Y Y'\to Y'$ 是 $f$ 的基变换。若 $f$ 射影，则 $f'$ 射影。
  \item 设 $f\colon X\to Y$ 与 $f'\colon X'\to Y'$ 是 $S$-概形之间的两个射影 $S$-态射。则 $f\times f'\colon X\times_S X'\to Y\times_S Y'$ 射影。
  \item 设 $f\colon X\to Y$ 与 $g\colon Y\to Z$ 是态射。若 $gf$ 射影且 $g$ 分离，则 $f$ 射影。
\end{enumerate}
\end{proposition}

\begin{proof}
我们只给出 (ii) 的证明，其余留给读者。设 $f\colon X\to Y$ 与 $g\colon Y\to S$ 是射影态射。它们可分解为复合
\[
X\to \bP^m_Y\to Y,
\]
\[
Y\to \bP^n_S\to S,
\]
使 $X\to \bP^m_Y$ 与 $Y\to \bP^n_S$ 是闭浸入。考虑以下图表
\begin{center}
\begin{tikzcd}
X \arrow[r] \arrow[dr] & \bP^m_Y \arrow[r] \arrow[d] & \bP^m_S\times_S\bP^n_S \arrow[d] \\
& Y \arrow[r] \arrow[dr] & \bP^n_S \arrow[d] \\
& & S.
\end{tikzcd}
\end{center}
注意上图中的方块是笛卡尔图，因为
\[
(\bP^m_S\times_S\bP^n_S)\times_{\bP^n_S}Y\cong \bP^m_S\times_S Y\cong \bP^m_Y.
\]
故复合 $X\to \bP^m_Y\to \bP^m_S\times_S\bP^n_S$ 是闭浸入。结合下面的 \ref{it:1.3.30}，这证明了 $gf$ 射影。
\end{proof}

\begin{lemma}\label{it:1.3.30}
存在闭浸入
\[
\bP^m_S\times_S\bP^n_S\to \bP^{(m+1)(n+1)-1}_S,
\]
它是 $S$-态射。
\end{lemma}

\begin{proof}
我们处理 $S=\Spec\bZ$ 的情形。一般情形可由基变换得到。有
\[
\bP^m_{\bZ}=\Proj \bZ[x_i]_{(i=0,\dots,m)},
\]
\[
\bP^n_{\bZ}=\Proj \bZ[y_j]_{(j=0,\dots,n)},
\]
\[
\bP^{(m+1)(n+1)-1}_{\bZ}=\Proj \bZ[z_{ij}]_{(i=0,\dots,m,\ j=0,\dots,n)}.
\]
对每个 $k\in\{0,\dots,m\}$ 与每个 $l\in\{0,\dots,n\}$，有满同态
\[
\bZ[z_{ij}]_{(z_{kl})}\to \bZ[x_i]_{(x_k)}\tensor_{\bZ}\bZ[y_j]_{(y_l)},\quad
\frac{z_{ij}}{z_{kl}}\mapsto \frac{x_i}{x_k}\tensor\frac{y_j}{y_l}.
\]
它诱导闭浸入
\[
D_+(x_k)\times D_+(y_l)\to D_+(z_{kl}).
\]
对任意 $k,k'\in\{0,\dots,m\}$ 与 $l,l'\in\{0,\dots,n\}$，有交换图表
\begin{center}
\begin{tikzcd}
\bZ[z_{ij}]_{(z_{kl})} \arrow[r] \arrow[d] & \bZ[x_i]_{(x_k)}\tensor_{\bZ}\bZ[y_j]_{(y_l)} \arrow[d] \\
\bZ[z_{ij}]_{(z_{kl}z_{k'l'})} \arrow[r] \arrow[u] & \bZ[x_i]_{(x_kx_{k'})}\tensor_{\bZ}\bZ[y_j]_{(y_ly_{l'})} \arrow[u] \\
\bZ[z_{ij}]_{(z_{k'l'})} \arrow[r] \arrow[u] & \bZ[x_i]_{(x_{k'})}\tensor_{\bZ}\bZ[y_j]_{(y_{l'})} \arrow[u]
\end{tikzcd}
\end{center}
其中中间的水平箭头由
\[
\frac{z_{ij}z_{i'j'}}{z_{kl}z_{k'l'}}\mapsto \frac{x_ix_{i'}}{x_kx_{k'}}\tensor\frac{y_jy_{j'}}{y_ly_{l'}}
\]
定义。故有交换图表
\begin{center}
\begin{tikzcd}
D_+(x_k)\times D_+(y_l) \arrow[r] \arrow[d] & D_+(z_{kl}) \arrow[d] \\
D_+(x_kx_{k'})\times D_+(y_ly_{l'}) \arrow[r] \arrow[u] & D_+(z_{kl}z_{k'l'}) \arrow[u] \\
D_+(x_{k'})\times D_+(y_{l'}) \arrow[r] \arrow[u] & D_+(z_{k'l'}). \arrow[u]
\end{tikzcd}
\end{center}
于是可以把闭浸入 $D_+(x_k)\times D_+(y_l)\to D_+(z_{kl})$ 粘起来，得到闭浸入 $\bP^m_{\bZ}\times\bP^n_{\bZ}\to \bP^{(m+1)(n+1)-1}_{\bZ}$。
\end{proof}

在 \ref{it:1.3.30} 的证明中构造的闭浸入 $\bP^m_S\times_S\bP^n_S\to \bP^{(m+1)(n+1)-1}_S$ 称为\emph{Segre 嵌入}\idx{segreyuqianru}{Segre 嵌入}。

\begin{tnote}
Segre 嵌入的靶空间维数 $(m+1)(n+1)-1$ 来自双指标 $z_{ij}$ 的个数：$i$ 有 $m+1$ 种取法，$j$ 有 $n+1$ 种取法，故 $\Proj$ 的坐标环有 $(m+1)(n+1)$ 个生成元。[依据：原书 p.56–57]
\end{tnote}

\begin{exercises}
\item \label{ex:1.3.1} 设 $X$ 是诺特拓扑空间，$\cF$ 是前层。假设对每个开子集 $U$ 与每个有限开覆盖 $(U_i)_{i\in I}$，以下条件成立：
\begin{enumerate}[label=(\roman*),leftmargin=2.8em]
  \item 若截面 $s,t\in \cF(U)$ 满足对每个 $i\in I$ 有 $s|_{U_i}=t|_{U_i}$，则 $s=t$。
  \item 设 $s_i\in \cF(U_i)$ 是截面，使对每对 $i,j\in I$ 有 $s_i|_{U_i\cap U_j}=s_j|_{U_i\cap U_j}$。则存在截面 $s\in \cF(U)$，使对每个 $i\in I$ 有 $s|_{U_i}=s_i$。
\end{enumerate}
证明 $\cF$ 是层。

\item \label{ex:1.3.2} 设 $(\cF_i,\phi_{ij})_{i\in I}$ 是 $X$ 上阿贝尔群层的有向系统，$\dlim_i\cF_i$ 是与前层 $U\mapsto \dlim_i\cF_i(U)$ 相伴的层。证明 $\dlim_i\cF_i$ 是 $(\cF_i,\phi_{ij})_{i\in I}$ 在层范畴中的有向极限，且对每个 $P\in X$ 有
\[
(\dlim_i \cF_i)_P\cong \dlim_i \cF_{i,P}.
\]
进一步假设 $X$ 是诺特拓扑空间。证明前层 $U\mapsto \dlim_i\cF_i(U)$ 是层。

\item \label{ex:1.3.3} 设 $(\cF_i,\phi_{ij})_{i\in I}$ 是 $X$ 上阿贝尔群层的逆向系统。证明前层 $U\mapsto \plim_i\cF_i(U)$ 是层，且它是 $(\cF_i,\phi_{ij})_{i\in I}$ 在层范畴中的逆向极限。

\item \label{ex:1.3.4} 设 $f\colon X\to Y$ 是概形的态射。假设存在 $Y$ 的开仿射覆盖 $V_i=\Spec B_i$（$i\in I$），使每个 $f^{-1}(V_i)$ 可被仿射开子概形 $U_{ij}=\Spec A_{ij}$（$j\in J_i$）覆盖，其中 $A_{ij}$ 是有限生成的 $B_i$-代数。设 $V=\Spec B$ 是 $Y$ 的仿射开子概形，$U=\Spec A$ 是 $X$ 的仿射开子概形，且 $f(U)\subset V$。
\begin{enumerate}[label=(\arabic*),leftmargin=2.4em]
  \item 证明 $f^{-1}(V)$ 可被仿射开子概形 $U_\lambda=\Spec A_\lambda$（$\lambda\in\Lambda$）覆盖，其中 $A_\lambda$ 是有限生成的 $B$-代数。
  \item 证明存在有限多个 $a_1,\dots,a_n\in A$，使 $D(a_1),\dots,D(a_n)$ 覆盖 $U$，且每个 $A_{a_i}$ 是有限生成的 $B$-代数。
  \item 证明 $A$ 是有限生成的 $B$-代数。
\end{enumerate}

\item \label{ex:1.3.5} 设 $S$ 是分次环，$\fa$ 是 $S$ 的齐次理想。证明典范同态 $S\to S/\fa$ 诱导闭浸入 $\Proj S/\fa\to \Proj S$。

\item \label{ex:1.3.6} 设 $\cC$ 是存在纤维积的范畴。考虑交换图表
\begin{center}
\begin{tikzcd}
X'' \arrow[r, "f'"] \arrow[d] & X' \arrow[r, "g'"] \arrow[d] & X \arrow[d] \\
S'' \arrow[r, "f"] & S' \arrow[r, "g"] & S,
\end{tikzcd}
\qquad
\begin{tikzcd}
X'' \arrow[r, "g'f'"] \arrow[d] & X \arrow[d] \\
S'' \arrow[r, "gf"] & S.
\end{tikzcd}
\end{center}
其中第一个图表的右方块为 (1)，左方块为 (2)；第二个图表中的方块为 (3)。假设方块 (1) 是笛卡尔图。证明方块 (2) 是笛卡尔图，当且仅当方块 (3) 是笛卡尔图。

\item \label{ex:1.3.7}
\begin{enumerate}[label=(\arabic*),leftmargin=2.4em]
  \item 设 $f\colon X\to S$ 与 $g\colon Y\to S$ 是集合的映射。证明它们的纤维积是 $(X\times_S Y,p,q)$，其中
  \[
  X\times_S Y=\{(x,y)\mid x\in X,\ y\in Y,\ f(x)=g(y)\},
  \]
  $p\colon X\times_S Y\to X$ 与 $q\colon X\times_S Y\to Y$ 分别是投影 $p(x,y)=x$ 与 $q(x,y)=y$。
  \item 利用 (1) 中纤维积的描述，对集合范畴证明 \ref{it:1.3.24}。
  \item 设 $\cC$ 是范畴，
  \begin{center}
  \begin{tikzcd}
  X' \arrow[r] \arrow[d] & X \arrow[d] \\
  S' \arrow[r] & S
  \end{tikzcd}
  \end{center}
  是 $\cC$ 中的交换图表。对每个对象 $Z\in\ob\cC$，它诱导集合范畴中的交换图表
  \begin{center}
  \begin{tikzcd}
  \Hom(Z,X') \arrow[r] \arrow[d] & \Hom(Z,X) \arrow[d] \\
  \Hom(Z,S') \arrow[r] & \Hom(Z,S).
  \end{tikzcd}
  \end{center}
  证明第一个图表在 $\cC$ 中是笛卡尔图，当且仅当对每个对象 $Z\in\ob\cC$，第二个图表在集合范畴中是笛卡尔图。
  \item 利用 (2) 与 (3)，对每个存在纤维积的范畴 $\cC$ 证明 \ref{it:1.3.24}。
\end{enumerate}

\item \label{ex:1.3.8} 集合
\[
I=\{(i_0,\dots,i_n)\mid i_0+\cdots+i_n=d,\ i_0,\dots,i_n\in \bZ_{\ge 0}\}
\]
有 $\binom{d+n}{n}$ 个元素。证明同态
\[
\bZ[y_{i_0\dots i_n}]_{(i_0,\dots,i_n)\in I}\to \bZ[x_0,\dots,x_n],\quad
y_{i_0\dots i_n}\mapsto x_0^{i_0}\cdots x_n^{i_n}
\]
诱导态射 $\bP^n_{\bZ}\to \bP^{\binom{d+n}{n}-1}_{\bZ}$。
\end{exercises}
